\documentclass[11pt]{article}
\usepackage[a4paper,margin=23mm]{geometry}
\usepackage{amsmath,amssymb,array,booktabs}
\usepackage[hidelinks]{hyperref}
\setlength{\parindent}{0pt}
\setlength{\parskip}{6pt}
\newcommand{\R}{\mathbb R}
\newcommand{\C}{\mathbb C}
\newcommand{\Sig}{\Sigma^*}
\newcommand{\dd}{\widetilde\partial}
\begin{document}
\begin{center}
{\Large No positive second free Stein discrepancy value}\\[4pt]
{\large Problem 00000007733: the fixed semicircle target}\\[3pt]
lizaixi01\qquad 25 October 2026
\end{center}

\textbf{The precise printed clause.}
The source ranges over all atomless distributions supported on $[-2,2]$,
asserts minimum $0$ at the semicircle law, and then asserts a
``next-smallest value.'' We use its usual meaning: a next \emph{distinct},
strictly positive value. The target is fixed throughout at the unit
semicircle, with potential $V_1(x)=x^2/2$. Support means containment in
$[-2,2]$; the source imposes no variance-one constraint.

We prove, from actual kernels and their full $L^2$ infimum, that for every
$c>0$ there is a genuine atomless operator law $\mu_s$ with
\[
 \operatorname{supp}(\mu_s)\subseteq[-2,2],\qquad
 0<\Sig(\mu_s\mid V_1)=1-s^2<c.
\]
Thus no positive second distinct value can exist in the printed class.
This addresses neither the atomic three-point assertion nor a Jacobi-family
classification. Exact-full-support, variance-normalized or stationary-value
repairs are outside this claim.

\textbf{The actual measures and bounded operators.}
Define the probability measure
\[
 d\nu(x)=\frac2\pi\sqrt{1-x^2}\,1_{[-1,1]}(x)\,dx,
 \qquad \mu_s=(x\mapsto2sx)_*\nu,\qquad 0<s\le1.
\]
Lean proves the density integral is $1$, atomlessness from absolute
continuity, and atomlessness of the pushforward from injectivity. The actual
almost-everywhere support bounds give the interval inclusion. The mean is
$0$ and the second moment is $s^2$, both computed as native measure
integrals. At $s=1$, change of variables identifies every continuous test
integral with the unit Wigner density
\[
 \int f\,d\mu_1=\int_{-2}^2 f(x)\frac{\sqrt{4-x^2}}{2\pi}\,dx.
\]
On the native complex Hilbert space $H_s=L^2(\mu_s)$, let $T_s f(x)=xf(x)$.
It is constructed as a continuous linear operator using the genuine
$L^\infty\times L^2\to L^2$ multiplication map. We prove
\[
 T_s=T_s^*,\quad\|T_s\|\le2,\quad
 \phi_s(B)=\langle1,B1\rangle,\quad\phi_s(1)=1,\quad
 \phi_s(B^*B)\ge0.
\]
For every $P\in\C[X]$, its native algebra evaluation satisfies
$\phi_s(P(T_s))=\int P(x)\,d\mu_s(x)$. The state is also proved tracial on
the actual commutative polynomial operator algebra; no trace assertion on
all bounded operators is used.

\newpage
\textbf{All polynomial tests, derived from the density.}
For a complex polynomial $P$, use the genuine divided difference
\[
 (\dd P)(x,y)=
 \begin{cases}P'(x),&x=y,\\[2pt]
 (P(x)-P(y))/(x-y),&x\ne y.
 \end{cases}
\]
We prove its diagonal derivative convention and measurability. Atomlessness
also gives that the diagonal has zero product measure.

Let $U_n$ and $T_n$ be the native Chebyshev polynomials. Actual weighted
integrals, using native $T_n$ orthogonality, give
\[
 \int U_n(x)\,d\nu(x)=\begin{cases}1,&n=0,\\0,&n>0.\end{cases}
\]
The algebraic convolution identity is proved in every degree:
\[
 2(x-y)\sum_{i=0}^n U_i(x)U_{n-i}(y)
   =U_{n+1}(x)-U_{n+1}(y).
\]
Its actual product integral and the coordinate--$U_n$ integral establish
the free integration identity on each $U_n$. The full complex polynomial
sequence is a native basis, so a proved submodule closure extends the result
to \emph{every} complex polynomial, not a finite collection of moments.
Under the actual dilation pushforward this becomes
\begin{equation}
 \int xP(x)\,d\mu_s(x)
   =s^2\int (\dd P)(x,y)\,d(\mu_s\otimes\mu_s)(x,y).
 \tag{1}
\end{equation}
All tests are proved to belong to every relevant native $L^p$ space.
The all-degree moment recursion is also explicitly derived by Fubini:
\[
 m_0=1,\quad m_1=0,\quad m_2=s^2,\qquad
 m_{n+2}=s^2\sum_{i=0}^n m_i m_{n-i},
 \quad m_n=\int x^n\,d\mu_s(x).
\]

\textbf{The full kernel domain and genuine infimum.}
We formalize Definition 1.8 of Fathi--Nelson in its one-variable measure
form from Section 4. Kernels range over the \emph{entire} native complex
Hilbert space $L^2(\mu_s\otimes\mu_s)$. Every $P\in\C[X]$ must satisfy the
Stein inner-product identity. Lean's conjugate-linear-first convention is
handled by proved integral/conjugation identities. No restriction to
constant, symmetric, polynomial, real or finite-dimensional kernels occurs.

Equation (1) proves that the actual constant function $A_s=s^2$ is an
admissible kernel. For any admissible kernel $A$, testing $P=X$ gives
$\langle1,A\rangle=s^2$. Since the constant-one vector has norm $1$,
Cauchy--Schwarz yields
\[
 1-s^2=|\langle1,A-1\rangle|\le\|A-1\|_2.
\]
The constant $A_s$ attains equality. The set of distances is explicitly
proved nonempty and bounded below before native \texttt{sInf} is used.
Consequently its actual full-domain infimum is
\[
 \boxed{\Sig(\mu_s\mid V_1)=\inf_{A\text{ Stein kernel}}\|A-1\|_2=1-s^2.}
\]

\newpage
\textbf{There is no next distinct value.}
At $s=1$ the same actual infimum is $0$. Given any $c>0$, put
\[
 \varepsilon=\min(c/4,1/4),\qquad s=1-\varepsilon.
\]
Then $0<s<1$ and
\[
 0<1-s^2=2\varepsilon-\varepsilon^2<c.
\]
The theorem \texttt{operator\_laws\_have\_no\_positive\_next\_value} includes
the real measure, its atomlessness/support, the actual self-adjoint bounded
operator, the normalized positive state, every complex polynomial law
identity and the strict discrepancy inequalities. The separate theorem
\texttt{no\_positive\_gap\_on\_admissible\_family} gives the logical
negation of a positive gap. Since this family lies inside the source's full
class, any proposed positive second distinct value for the full class has
a smaller positive value inside this family.

\textbf{Formalization boundary.}
The one-variable measure/product-measure formulation is the published
definition in Fathi--Nelson, Section 4, read together with the all-kernel
infimum of Definition 1.8. The proof constructs those actual measures and
their native $L^2$ spaces directly. It also constructs their coordinate
operators and proves their complete polynomial distribution identities.
A general von Neumann algebra/GNS/tensor-product reconstruction and a
general spectral functional-calculus reconstruction are not repeated in
Lean. No added axiom substitutes for either reconstruction. No identity
between Stein discrepancy and free Fisher information is claimed.

\textbf{Proof map.}
\begin{center}
\begin{tabular}{>{\raggedright\arraybackslash}p{.38\linewidth}>{\raggedright\arraybackslash}p{.55\linewidth}}
\toprule
Actual construction & Main declarations\\
\midrule
Density, probability and all basis tests &
 \texttt{baseLaw}, \texttt{base\_U\_integral}, \texttt{all\_base\_tests}\\
Pushforward and all moments &
 \texttt{law}, \texttt{law\_free\_integration\_by\_parts},
 \texttt{all\_moments\_recurrence}\\
All-kernel infimum and attainment &
 \texttt{IsSteinKernel}, \texttt{kernel\_norm\_lower},
 \texttt{discrepancy\_eq}\\
Bounded self-adjoint operator law &
 \texttt{coordinateOperator}\newline \texttt{coordinateOperator\_distribution}\\
\bottomrule
\end{tabular}
\end{center}

\textbf{Reproduction and coverage.}
The complete project pins Lean 4.33.0 and Mathlib commit\\
\texttt{db584cd6d46c92f209a44c0f1c829460d327499d}.
Check the runtime with \texttt{lean --version}. From the package directory,
\texttt{bash lean/verify.sh} runs:
\begin{quote}\small\ttfamily
lake build +Main\\
lake env lean -DwarningAsError=true Check.lean\\
lake env leanchecker --fresh Main
\end{quote}
The dynamic check enumerates every owned namespace declaration, including
generated declarations, in addition to the complete explicit-name list.
Only \texttt{propext}, \texttt{Classical.choice}, and \texttt{Quot.sound} are
allowed. Production timings are diagnostic and establish no controlled
research-efficiency improvement.

\textbf{Primary reference.}
M. Fathi and B. Nelson, \emph{Free Stein kernels and an improvement of the
free logarithmic Sobolev inequality}, arXiv:1605.05828v1, 19 May 2016:
printed p.4 (difference quotients), p.7 (Definition 1.8), and pp.13--14
(the fixed unit-semicircle density and one-variable measure formula).
\url{https://arxiv.org/abs/1605.05828v1}.
\end{document}
